La figura mostra l'esquema d'un relé. Quan circula un corrent elèctric per la bobina, l'extrem inferior de l'imant (nord) és atret per la bobina i el moviment es transmet per un pivot, de manera que es tanca el circuit B.

\figura[0.5]{fig1.png}

\begin{apartat}{1}
Especifiqueu clarament quin ha de ser el sentit del corrent elèctric a la bobina perquè s'activi el relé (i es tanqui el circuit B) i dibuixeu les línies del camp magnètic generat per la bobina en aquesta situació.
\end{apartat}

\begin{apartat}{1}
En unes proves observem que el mecanisme no fa prou força per a tancar el contacte. Indiqueu quin efecte tindria sobre el dispositiu cadascuna de les modificacions següents:
\begin{enumerate}
  \item[1)] Augmentar la intensitat del corrent que circula per la bobina.
  \item[2)] Situar un material ferromagnètic al nucli de la bobina.
  \item[3)] Fer passar per la bobina un corrent altern en comptes d'un corrent continu.
\end{enumerate}
\end{apartat}
